\SourceLeaf{354}%
sich von Stelle zu Stelle auf den Flächen ändert. Bei der \textsc{Mercator}-Projektion ist die Vergrößerungszahl zum Beispiel umso größer, je weiter die betreffenden Stellen vom Äquator entfernt sind. Bei der stereografischen Projektion ist sie umso größer, je weiter die Stellen vom Augenpunkt entfernt sind.

Jede gegebene Fläche lässt sich auf unendlich viele Arten konform auf eine andere gegebene Fläche abbilden. Die allgemeine Lösung umfasst sie alle, indem sie eine beliebige Funktion enthält, die frei oder entsprechend dem jeweiligen Zweck gewählt werden kann. Soll nur ein Teil der einen Fläche abgebildet werden, wählt man in der Regel am besten eine konforme Abbildung, bei der die Schwankungen des Vergrößerungsverhältnisses innerhalb des abzubildenden Flächenteils in möglichst engen Grenzen bleiben.

Die Aufgabe, die Ellipsoidoberfläche konform auf die Kugeloberfläche abzubilden, wird in der angeführten Schrift unter den Beispielen eigens behandelt. Der allgemeinen Lösung sind zwei besondere Lösungen beigefügt. Die eine eignet sich vor allem für die Abbildung der gesamten Ellipsoidoberfläche. Die andere ist weit zweckmäßiger, wenn nur ein mäßig großer Teil der als ellipsoidisch betrachteten Erdoberfläche konform auf eine Kugeloberfläche abgebildet werden soll, wie es bei konkreten Anwendungen in der Geodäsie stets der Fall ist. Wie gut diese zweite Abbildungsart die oben genannte Anforderung erfüllt, lässt sich einem dort angeführten Beispiel entnehmen. Innerhalb einer Zone von fünf Breitengraden beträgt die Schwankung des Vergrößerungsverhältnisses nur \(\frac1{530000}\). Dort ist ferner in groben Zügen angedeutet, wie sich eine konforme Abbildung überhaupt zur Berechnung eines Dreiecksnetzes verwenden lässt. Die eigentliche Ausführung und die Anwendung auf diese bestimmte Abbildungsart blieben einer späteren Bearbeitung vorbehalten.

Die vorliegende Abhandlung soll nun diese Verpflichtung einlösen, wenn auch etwas anders als ursprünglich vorgesehen. Ihr liegt eine dritte besondere Lösung der Aufgabe zugrunde, die sich von der eben erwähnten unterscheidet und den angestrebten Zweck noch vollkommener erreicht. Hier müssen wir uns damit begnügen, nur eine allgemeine Vorstellung davon zu geben.
